%%%%%%%%%%%%%%%%%%%%%%%%%%%%%%%%%%%%%%%%%
% Medium Length Professional CV
% LaTeX Template
% Version 2.0 (8/5/13)
%
% This template has been downloaded from:
% http://www.LaTeXTemplates.com
%
% Original author:
% Rishi Shah
%
% Important note:
% This template requires the resume.cls file to be in the same directory as the
% .tex file. The resume.cls file provides the resume style used for structuring the
% document.
%
%%%%%%%%%%%%%%%%%%%%%%%%%%%%%%%%%%%%%%%%%

%----------------------------------------------------------------------------------------
%	PACKAGES AND OTHER DOCUMENT CONFIGURATIONS
%----------------------------------------------------------------------------------------

\documentclass{resume} % Use the custom resume.cls style
\usepackage[utf8]{inputenc}
\usepackage{hyperref}

\usepackage[left=0.45in,top=0.6in,right=0.35in,bottom=0.6in]{geometry} % Document margins
\newcommand{\tab}[1]{\hspace{.2667\textwidth}\rlap{#1}}
\newcommand{\itab}[1]{\hspace{0em}\rlap{#1}}
\name{\large{João Francisco Pugliese} \\ } % Your name

\begin{document}

%----------------------------------------------------------------------------------------
%	EDUCATION SECTION
%----------------------------------------------------------------------------------------


220 College Ave. \hfill Email: jfpugli@stanford.edu \\
Palo Alto, California \hfill Date of birth: February 6, 1995 \\


%----------------------------------------------------------------------------------------
%	EDUCATION SECTION
%----------------------------------------------------------------------------------------

\begin{rSection}{Education}

{\bf Ph.D. Economics} Stanford University \hfill {\em exp. 2027}
\begin{tabular}{ll}
    Advisors: Prof. Matthew Gentzkow and Prof. Arun Chandrasekhar
\end{tabular}

{\bf M.A. Economics} São Paulo School of Economics (FGV-EESP) \hfill {\em 2020} \\
\begin{tabular}{ll}
    Advisor: Prof. Bruno Ferman  \\
    Thesis Title: Reference-dependence and labor supply decisions.
\end{tabular} \\
\\{\bf B.A. Economics}, São Paulo School of Economics (FGV-EESP) \hfill {\em 2017}
\\
\begin{tabular}{l}
    Advisor: Prof. Leonardo Weller  \\
    Thesis Title: Human capital persistence in the Jesuit missions of southern Brazil.
\end{tabular} \\

\end{rSection}

%----------------------------------------------------------------------------------------
%	References
%----------------------------------------------------------------------------------------

\begin{rSection}{References}

\begin{tabular}{@{}p{0.305\textwidth}p{0.305\textwidth}p{0.305\textwidth}@{}}
{\bf Arun Chandrasekhar} (co-primary) & {\bf Matthew Gentzkow} (co-primary) & {\bf Melanie Morten} \\
Department of Economics & Department of Economics & Department of Economics \\
Stanford University & Stanford University & Stanford University \\
\href{mailto:arungc@stanford.edu}{arungc@stanford.edu} & \href{mailto:gentzkow@stanford.edu}{gentzkow@stanford.edu} & \href{mailto:memorten@stanford.edu}{memorten@stanford.edu} \\
\end{tabular}

\end{rSection}
%--------------------------------------------------------------------------------
%   Research Papers
%-----------------------------------------------------------------------------------------------
\begin{rSection}{Research Papers}

{\bf Consumer misperceptions and redistribution in financial products: evidence \\ from Brazilian consórcios (JMP)} \hfill {\em October 2026} \\
{\it with Otávio Braga and Lucas I. Teixeira} \\

{\bf Inflation beliefs and the indexation of household debt} \hfill {\em October 2026} \\
{\it with Otávio Braga and Lucas I. Teixeira} \\

{\bf Good news and bad news: the effects of consuming professional news outlets} \hfill {\em September 2026} \\
{\bf During the 2022 Brazilian Elections} \\
{\it with Caio Castro} \\

{\bf \href{https://joaofpugliese.github.io/assets/pdf/CDKPRS.pdf}{Blue Spoons: sparking communication about appropriate technology use}} \hfill {\em August 2022} \\
{\it with Arun G. Chandrasekhar, Esther Duflo, Michael Kremer, Jonathan Robinson, and Frank Schilbach} \\


{\bf \href{https://joaofpugliese.github.io/assets/pdf/BP.pdf}{Can income targeting be explained by dynamic sample selection?}} \hfill {\em July 2021} \\
{\it with Pedro Bessone} \\

\end{rSection}

%----------------------------------------------------------------------------------------
%	Academic positions
%----------------------------------------------------------------------------------------

\begin{rSection}{Academic Experience}
{\bf Research Assistant} \hfill {\em 2021-2022} \\
{\it Description:} I worked as Research Assistant for professor Melanie Morten on projects relating to transmigration, with a specific focus on Northern Triangle (El Salvador, Guatemala and Honduras) citizens migrating to the United States. \\ \\
{\bf Research Assistant} \hfill {\em 2019-2020} \\
{\it Description:} I worked as Research Assistant for professors Gautam Rao, Frank Schilbach and Heather Schofield in the paper  \href{https://economics.mit.edu/sites/default/files/2022-09/economic-sleep-qje.pdf}{``The Economic Consequences of
Increasing Sleep Among the Urban Poor"}. I conducted econometric and data analysis with large datasets.
\end{rSection}

%----------------------------------------------------------------------------------------
%	Professional Experience
%----------------------------------------------------------------------------------------

\begin{rSection}{Professional Experience}

\smallskip
{\bf Uber Applied Sciences Ph.D. Summer Internship} \hfill {\em 2025} \\
{\it Description:} At Uber, I spent 12 weeks investigating the role of spillovers in matching markets and demand modeling, combining large scale data from experiments to enhance algorithmic pricing. I produced two final reports focusing on distinct parts of the company's demand estimation process.

\end{rSection}
%----------------------------------------------------------------------------------------
%	Teaching Experience
%----------------------------------------------------------------------------------------

\begin{rSection}{Teaching Experience}

{\bf TA Media Markets and Social Good ECON 47} (undergrad.) at Stanford (Prof. Matt Gentzkow) \hfill {\em 2025} \\
\smallskip
{\bf TA Game Theory ECON 51} (undergrad.) at Stanford (Prof. Chris Makler) \hfill {\em 2024, 2025} \\
\smallskip
{\bf TA Political Economy I} (grad.) at Stanford (Prof. Matt Gentzkow and Andy Hall) \hfill {\em 2025} \\
\smallskip
{\bf TA Probability and Statistics ECON 102A} (undergrad.) at Stanford (Prof. Scott McKeon) \hfill {\em 2024} \\
\smallskip
{\bf TA Econometrics  ECON 102B} (undergrad.) at Stanford (Prof. Scott McKeon) \hfill {\em 2023, 2025} \\
\smallskip
{\bf TA Data Science ECON 108} (undergrad.) at Stanford (Prof. Han Hong) \hfill {\em 2020} \\
\smallskip
{\bf TA Stats Camp} (grad.) at FGV-EESP (Prof. Priscilla Albuquerque) \hfill {\em 2020} \\
\smallskip
{\bf TA Econometrics III} (undergrad.) at FGV-EESP (Prof. Priscilla Albuquerque) \hfill {\em 2020} \\
\smallskip
{\bf TA Math Camp (Analysis)} (grad.) at FGV-EESP (Prof. Victor Filipe Martins-da-Rocha) \hfill {\em 2019} \\
\smallskip
{\bf TA Econometrics I} (undergrad.) at FGV-EESP (Prof. Priscilla Albuquerque) \hfill {\em 2019} \\
\smallskip

\end{rSection}

%--------------------------------------------------------------------------------
%    Awards and Grants
%-----------------------------------------------------------------------------------------------
\begin{rSection}{Grants and Awards}
{\bf Weiss Fund Pilot Award} \hfill {\em 2026} \\
{\bf Department of Economics Teaching Assistant Award (3x)} \hfill {\em 2023, 2024, 2026} \\
{\bf SIEPR Dissertation Fellowship} \hfill {\em 2025-2026} \\
{\bf Centennial Teaching Assistant (CTA) Award} \hfill {\em 2025} \\
{\bf SIEPR George P. Schultz Grant} \hfill {\em 2025} \\
{\bf Stanford's Graduate Research Opportunties (GRO) Grant} \hfill {\em 2025} \\
{\bf King Center Advanced Research Grant} \hfill {\em 2022, 2025} \\
{\bf Scholarship from CAPES to fund Master's Program} \hfill {\em 2018-2020} \\
\end{rSection}

%----------------------------------------------------------------------------------------
%	Programming
%----------------------------------------------------------------------------------------

\begin{rSection}{Programming}

Stata (proficient), \LaTeX, R (proficient), Python (intermediate)
\end{rSection}

%----------------------------------------------------------------------------------------
%	Languages
%----------------------------------------------------------------------------------------

\begin{rSection}{Languages}

Portuguese (native), English (fluent)
\end{rSection}


\end{document}
